\documentclass[11pt]{article}
\usepackage[margin=1in]{geometry}
\usepackage{amsmath,amssymb}
\usepackage{tikz}
\usepackage{float}
\usepackage{graphicx}
\usepackage[hidelinks]{hyperref}

\setlength{\parindent}{0pt}
\setlength{\parskip}{0.7\baselineskip}

\begin{document}

\section{Game Search}

This is squarely in the `other solvers' category, but I've had a few requests.

What do we mean by game search?

\begin{quote}
\emph{(tic-tac-toe example)}
\end{quote}

\textbf{Minimax} is the basic game search. The name comes from finding the \emph{min} of the \emph{max}.

(courtesy of Wikipedia, some pseudocode)

\begin{verbatim}
function minimax( node, depth, maximizingPlayer ) is
    if depth = 0 or node is a terminal node then
        return the heuristic value of node
    if maximizingPlayer then
        value := -infinity
        for each child of node do
            value := max( value, minimax( child, depth - 1, FALSE ) )
        return value
    else (* minimizing player *)
        value := +infinity
        for each child of node do
            value := min( value, minimax( child, depth - 1, TRUE ) )
        return value
(* Initial call *)
minimax( origin, depth, TRUE )
\end{verbatim}

(also from Wikipedia, an example tree for a 2-player 2-option game with example leaf values)

\begin{figure}[H]
    \centering
    \includegraphics[width=0.75\textwidth]{image-1.png}
    \caption{An example of a minimax tree for a two-player game. \emph{(By Nuno Nogueira (Nmnogueira) -- own work using \url{http://en.wikipedia.org/wiki/Image:Minimax.svg}, created in Inkscape by author, CC BY-SA 2.5, \url{https://commons.wikimedia.org/w/index.php?curid=2276653}.)}}
\end{figure}

Let's construct this tree with values for our tic-tac-toe example.

What does this tell us about games? How can we interpret the value at the root?

\subsection{How can we make this faster?}

Without improvements, minimax will construct the entire game tree.

We can deploy many tricks!
\begin{itemize}
    \item alpha-beta pruning
    \item hash tables for repeat states
    \item node ordering (maybe find a win faster)
    \item trickery to make at-node computation faster
    \item if we're not likely to solve exactly, we can use node value heuristics
\end{itemize}

Two now-old things that were exciting at the time:
\begin{itemize}
    \item Alpha-Go \url{https://en.wikipedia.org/wiki/AlphaGo}, via Monte Carlo Tree Search
    \item Monte Carlo Tree Search, more generally
\end{itemize}

\begin{figure}[H]
    \centering
    \includegraphics[width=0.9\textwidth]{image.png}
    \caption{An example showing Monte Carlo Tree Search. \emph{(By Robert Moss -- own work, CC BY-SA 4.0, \url{https://commons.wikimedia.org/w/index.php?curid=111182752}.)}}
\end{figure}

\end{document}
